\section{引言}

2021年，国务院印发《2030年前碳达峰行动方案》，提出推动运输工具装备低碳转型，加快城市绿色货运配送发展\cite{ref:55}。城市配送连接仓储节点与终端客户，车辆运行频繁，服务时段集中，燃料消耗和碳排放直接影响物流企业的运营成本与减排目标。在低碳转型的政策要求和降低运营成本的经营需要双重推动下，优化配送车队的调度方案与能耗管理，已成为运筹学与物流管理领域的重要研究课题。

在燃油车尚不能完全退出、电动车持续投入运营的过渡阶段，由两类车辆共同承担配送任务已成为车队更新过程中具有现实意义的组织方式。燃油车续驶里程充裕、补能便捷，电动车运行能耗较低且无尾气排放，但两类车辆在固定成本、载重能力、续驶里程和补能方式等方面存在显著差异，车型配置与路径规划需要统一确定。李得成等\cite{ref:6}建立电动车与燃油车的混合车辆路径模型，设计分支定价算法联合优化车型配置与配送路径。陈婉茹等\cite{ref:23}在多配送中心混合车队中考虑车辆速度对能耗的影响，采用机械功率模型分别计算两类车辆的实际能耗，并将碳交易成本纳入优化目标。这些研究表明，混合车队的运营特征只有在车型选择、路径安排和能耗核算统一建模时才能得到准确刻画。

从理论视角来看，混合车队配送属于车辆路径问题（Vehicle Routing Problem，VRP）的重要扩展。VRP是组合优化领域的经典问题，自提出以来在配送调度、城市物流和供应链管理等领域得到广泛应用。随着城市配送组织方式和能源结构不断变化，传统的单车场、同质车队和静态需求假设已难以刻画现实运营中的多种限制，研究逐步扩展到多车场、多趟配送、电动车补能以及动态需求等情形。上述扩展分别描述配送系统中某一类关键限制，也为研究多种因素共同作用下的配送优化奠定了理论基础。

车辆能耗和补能方式直接影响混合车队路径的可行性与成本结构。Goeke和Schneider\cite{ref:27}建立油电混合车队路径模型，将车辆载重和路段坡度等因素纳入能耗计算，更准确地反映不同车型在同一路段上的能耗差异。在电动车充电方面，Keskin和Çatay\cite{ref:45}提出部分充电策略，允许电动车根据后续路径需要灵活选择补电量。Montoya等\cite{ref:montoya}采用分段线性函数刻画电池的非线性充电过程，反映充入相同电量所需时间随荷电状态递增的物理特性。Froger等\cite{ref:12}进一步考虑充电站的容量限制，研究多辆电动车在有限充电设施下的调度问题。当电价随时段变化时，车辆的到站时刻和充电区间将改变充电费用，Deng等\cite{ref:deng}分析不同充电定价方案对物流电动车队运营成本的影响。Shi等\cite{ref:93}在混合车队双目标路径模型中引入分时电价，研究电价时段对充电决策和车辆调度的作用。此外，电网碳排放强度在一天中随发电结构变化而波动，Li等\cite{ref:52}的研究表明，合理安排充电时段可以降低电动车充电环节的碳排放。将非线性充电过程与分时电价和时变碳强度在同一时间尺度上协调考虑，是混合车队能耗与排放建模尚待深入的方向。

当配送网络由单一车场扩展为多个车场时，不同车场的客户、车辆和充电设施可以协同使用，扩大各车场的服务范围。Fernández等\cite{ref:49}提出共享客户的协作配送模型，允许各车场交叉服务对方区域内的客户以降低总成本。Wang等\cite{ref:7}在多车场电动车路径问题中研究充电站的共享使用，将电动车充电资源纳入多车场协作框架。与多车场协作相伴随，多趟配送允许同一车辆在返回车场后继续承担后续任务，在车辆数量有限时提高车队利用率。Zhen等\cite{ref:44}研究带时间窗和释放时刻约束的多车场多趟路径问题。Şahin和Yaman\cite{ref:sahin}针对异质车队的多车场多趟问题设计分支定价算法进行精确求解。多车场协作所形成的成本节约需要在成员之间合理分配，以维持各方参与意愿。Soriano等\cite{ref:8}研究多车场路径问题中的利润公平分配。饶卫振等\cite{ref:91}考虑企业间服务质量差异，研究协作配送中的成本分摊方法。多车场与多趟运营结合后，同一车辆在连续趟次间的返场、补货、补电和再次出发必须保持时间与资源的连续，协作产生的成本节约也需在各方之间依据实际贡献加以分配。

以上研究多以需求信息在规划开始前完全已知为前提，而城市配送中的订单往往在运营过程中持续到达或发生变化。Psaraftis\cite{ref:14}较早提出动态车辆路径问题的研究框架，此后该领域围绕信息更新方式与重规划策略形成了丰富的研究。李阳等\cite{ref:18}研究动态需求下车辆路径问题的周期性优化，通过设定重规划时刻将动态问题分解为一系列静态子问题。姜广田等\cite{ref:51}研究动态需求下的多车型路径优化问题。邱莹莹\cite{ref:71}针对低碳混合车队配送，采用定时与定量联合触发策略应对动态到达的订单。对于多车场多趟混合车队，每次重规划不仅需要更新待服务的客户集合，还需保留各车辆已执行任务的位置、时刻、载重和电量状态，使前后方案在时间与资源上保持连续可行。

为直观反映相关研究对多车场、混合车队、多趟配送、动态需求、分时充电排放和成本分摊等问题特征的考虑情况，选取与本文关系较为密切的研究进行比较，结果见表\ref{tab:literature-comparison}。

\begin{table}[H]
\centering
\caption{相关研究内容比较}
\label{tab:literature-comparison}
\setptabsetup
\begin{tabular*}{\linewidth}{@{\extracolsep{\fill}}lcccccc@{}}
\toprule
文献 & 多车场 & 混合车队 & 多趟配送 & 动态需求 & 分时充电排放 & 成本分摊\\
\midrule
Goeke和Schneider\cite{ref:27} & --- & $\checkmark$ & --- & --- & --- & ---\\
Zhen等\cite{ref:44} & $\checkmark$ & --- & $\checkmark$ & --- & --- & ---\\
陈婉茹等\cite{ref:23} & $\checkmark$ & $\checkmark$ & --- & --- & --- & ---\\
Wang等\cite{ref:7} & $\checkmark$ & --- & --- & --- & --- & ---\\
Soriano等\cite{ref:8} & $\checkmark$ & --- & --- & --- & --- & $\checkmark$\\
邱莹莹\cite{ref:71} & --- & $\checkmark$ & --- & $\checkmark$ & --- & ---\\
姜广田等\cite{ref:51} & --- & --- & --- & $\checkmark$ & --- & ---\\
Li等\cite{ref:52} & --- & --- & --- & --- & $\checkmark$ & ---\\
本文 & $\checkmark$ & $\checkmark$ & $\checkmark$ & $\checkmark$ & $\checkmark$ & $\checkmark$\\
\bottomrule
\end{tabular*}
\end{table}

由表\ref{tab:literature-comparison}可见，既有研究通常围绕其中一类或数类问题特征展开，多种运营特征在同一配送系统中的协同作用仍有进一步研究空间。

综合来看，现有研究已分别在混合车队能耗建模、非线性充电、分时电价、多车场协作与公平分配、多趟配送以及动态需求等方面取得了成果，但当多种因素同时出现时，仍存在以下不足：
1）非线性充电过程与分时电价、时变电网碳强度之间缺少统一的时段对应关系，现有研究或单独考虑非线性充电特性，或单独引入分时电价，尚未将充电过程中电量变化、费用核算与碳排放计算在同一时间尺度上加以协调；
2）多车场协作与多趟配送共同出现时，同一车辆在连续趟次间的补货、补电与出发时刻及载重电量约束必须保持连续，现有模型较少同时处理多车场资源共享与同一车辆的趟次衔接约束；
3）动态需求下的重规划研究多针对单车场或同质车队，对于多车场多趟混合车队如何在保留已执行任务的基础上同步更新客户集合与车辆状态，并保持成本与排放的核算口径一致，尚缺少系统的建模方法。

基于此，本文研究时变电网碳强度下多车场动态协同混合车队路径优化问题，建立以车辆启动、行驶、充电、油耗和碳排放五类成本之和最小为目标的优化模型，设计混合遗传搜索算法进行求解。论文主要工作如下：
1）综合考虑多车场协作、燃油车与电动车混合车队、多趟配送、非线性充电和动态需求，建立车辆路径与趟次衔接、充电安排相协调的配送优化模型；
2）依据实际充电区间与电价时段、电网碳强度时段的对应关系，分别核算充电费用和电动车间接碳排放；
3）针对多车场、多趟、充电和动态需求等问题特征，设计相应的解表示、邻域搜索和动态重规划方法；
4）通过公开算例检验算法的求解能力，并分析车队动力配置、充电策略、协作配送、成本分配和动态需求对配送方案的影响。
